\documentclass[11pt]{article}

\usepackage[margin=1in]{geometry}
\usepackage[T1]{fontenc}
\usepackage[utf8]{inputenc}
\usepackage{lmodern}
\usepackage{amsmath, amssymb, amsfonts}
\usepackage{booktabs}
\usepackage{array}
\usepackage{tabularx}
\usepackage{xcolor}
\usepackage{enumitem}
\usepackage{hyperref}

\hypersetup{colorlinks=true,linkcolor=blue!50!black,citecolor=blue!50!black,urlcolor=blue!50!black}

\newcommand{\AWR}{\mathrm{AWR}}
\newcommand{\ESS}{\mathrm{ESS}}
\newcommand{\FR}{\mathrm{FR}}
\newcommand{\EI}{\mathrm{EI}}
\newcommand{\Rdist}{R_{\mathrm{dist}}}

\title{A Coherent Extractability Story for Offline GCRL}
\author{}
\date{}

\begin{document}
\maketitle

\section{First principles: what AWR actually sees}

Offline goal-conditioned RL methods learn different objects: Bellman values, quasimetric distances,
contrastive reachability scores, or value-progress functions. However, under a shared AWR extractor,
the actor does not directly see these objects. It sees only an advantage value
\[ A_{\mathrm{alg}}(s,a,g) \]
that is converted into a behavioral cloning weight
\[ w_i = \min\left\{ \exp\left(\alpha A_{\mathrm{alg}}(s_i,a_i,g_i)\right), w_{\max} \right\}. \]

Thus, an advantage signal is useful for extraction only if it satisfies three conditions:
\begin{enumerate}[leftmargin=2em]
    \item \textbf{Correct ranking:} useful dataset actions should receive higher advantage than irrelevant actions.
    \item \textbf{Useful scale:} the advantage gaps should be large enough to shape behavior, but not so large that AWR collapses onto a tiny unstable tail.
    \item \textbf{Reliable support:} high-weight samples should correspond to behavior-supported actions that can actually be imitated from the dataset.
\end{enumerate}

The learning rate $\eta$ affects the learned critic, distance, or representation.
The AWR temperature $\alpha$ controls how aggressively the learned signal is converted into policy-imitation weights.
The central object of the paper is therefore the response surface in the $(\eta,\alpha)$ plane.

\section{Why progress in AntMaze is not the same as progress in Cube}

\subsection{AntMaze: local transition progress usually aligns with goal progress}

In AntMaze, the goal is essentially a two-dimensional navigation target. A simplified state can be written as
$s_t = (x_t,y_t,\text{proprioception})$, $g=(x_g,y_g)$.
A useful transition often moves the agent along a path-like structure toward the goal: $s_t \rightarrow s_{t+1}$.
If $s_{t+1}$ is closer to the goal along the maze path, then for a distance-like value $d(s_t,g) > d(s_{t+1},g)$,
and for a value-like function, $V(s_{t+1},g) - V(s_t,g) > 0$.

\subsection{Cube: useful manipulation may not look like immediate goal progress}

In Cube, the goal is object position $g = p_{\mathrm{cube}}^{\star} \in \mathbb{R}^3$, but the state contains robot proprioception and gripper states.
During a useful transition (e.g., reaching for the cube), the object may not move: $\|p_{\mathrm{cube},t+1}-p_{\mathrm{cube}}^\star\| \not< \|p_{\mathrm{cube},t}-p_{\mathrm{cube}}^\star\|$.
This changes the meaning of a progress signal: $V(s_{t+1},g)-V(s_t,g)$ may remain useful if it captures delayed success, but geometric signals become less aligned.

\section{Actor-facing advantage semantics}

The four algorithms expose different quantities to AWR.

\begin{table}[htb]
\centering
\small
\renewcommand{\arraystretch}{1.2}
\begin{tabular}{lll}
\toprule
Algorithm & Actor-facing advantage & Informal semantics \\
\midrule
GCIQL & $Q(s,a,g)-V(s,g)$ & Is this action better than the dataset-average? \\
GCIVL & $V(s',g)-V(s,g)$ & Did the transition make local value progress? \\
QRL   & $d_\phi(s,g)-d_\phi(s',g)$ & Did the transition reduce quasimetric distance? \\
CRL   & $f_\theta(s,a,g)-f^V_\theta(s,g)$ & Did this action increase reachability? \\
\bottomrule
\end{tabular}
\caption{The shared AWR extractor receives semantically different advantage signals across algorithms.}
\end{table}

\section{Diagnostics measured}

\subsection{Semantic reliability diagnostics}

These diagnostics ask: \textit{does the advantage function know which goal is the right one?}
They are computed from a $256\times256$ cross-goal advantage matrix $A(s_i, g_j)$, where the
diagonal ($i=j$) contains the matched future goal and the off-diagonal entries contain random goals.

\begin{itemize}[leftmargin=2em, itemsep=4pt]

  \item \textbf{FR-AUC} (\emph{future-random area under the ROC curve}).
  Pick one matched pair $(s_i, g_i^+)$ and one random pair $(s_i, g_j^-)$.
  FR-AUC is the probability that the matched pair gets a higher advantage score.
  $0.5$ means the critic cannot tell the right goal from a wrong one (coin flip).
  Above $0.5$ means it ranks the correct goal higher on average.
  Below $0.5$ means it is \emph{inverted}: it actively prefers random goals over the true one.

  \item \textbf{Gap mean} ($\overline{\delta}$).
  The average margin by which the matched goal wins: $\mathbb{E}[A^+ - A^-]$.
  FR-AUC only counts wins and losses; gap mean measures the \emph{size} of the win.
  Because AWR exponentiates advantages, a large gap creates concentrated weights even if
  FR-AUC is only moderately above $0.5$. Negative gap mean means the signal is inverted.

  \item \textbf{Extractability Index (EI)}.
  Gap mean divided by the standard deviation of the gaps: $\overline{\delta}/\sigma_\delta$.
  This is a signal-to-noise ratio. A large average margin that is \emph{consistent} across all pairs
  scores high. A large margin that varies wildly scores low. High EI means the signal is
  both strong and reliable enough to exploit.

  \item \textbf{MRR} (\emph{mean reciprocal rank}).
  For each observation $s_i$, rank all $256$ goals by their advantage score.
  MRR is the average of $1/\mathrm{rank}$ of the true goal.
  MRR$=1$ means the true goal is always ranked first; MRR$\approx 1/256$ means it is essentially random.
  Unlike FR-AUC, MRR is sensitive to whether the true goal reaches the very top, not just
  whether it beats one random competitor.

  \item \textbf{Prec@1 / Prec@5}.
  Fraction of observations where the true goal is ranked \#1 (or in the top 5) out of 256.
  Stricter than MRR: does the critic consistently put the right goal at the very top of the ranking?

  \item \textbf{$A^+$ mean and $A^-$ std}.
  The absolute scale of matched-goal advantages and the spread of random-goal advantages.
  Together they indicate whether the two distributions are well-separated or heavily overlapping.

\end{itemize}

\subsection{AWR extraction-conditioning diagnostics}

These diagnostics ask: \textit{given the advantage signal, does AWR produce a useful actor update?}
They use only the matched-goal (diagonal) advantages $A_i^+$ with the algorithm's configured $\alpha$
to compute clipped weights $w_i = \min\{\exp(\alpha A_i^+),\, w_{\max}\}$.

\begin{itemize}[leftmargin=2em, itemsep=4pt]

  \item \textbf{ESS} (\emph{effective sample size}).
  $\ESS = (\sum_i w_i)^2 / (N \sum_i w_i^2)$, normalized to $[0,1]$.
  Measures how many transitions the actor update effectively uses.
  $\ESS = 1$ means all samples have equal weight (fully diffuse cloning).
  $\ESS \to 0$ means one sample dominates (the update collapses onto a single transition).
  Directly measures how concentrated the behavioral cloning step is.

  \item \textbf{Saturation mass}.
  Fraction of total weight held by samples that hit the clip ceiling ($w_i = w_{\max}$).
  If this is high, most of the actor update is driven by samples whose exact advantage no longer
  matters---only that they exceeded the ceiling. The update becomes insensitive to the advantage
  values of the most influential transitions.

  \item \textbf{Fraction clipped}.
  Fraction of \emph{samples} (not weight) that hit the ceiling.
  Unlike saturation mass, this counts how many transitions are clipped, not how much weight they carry.
  Low fraction clipped with low ESS means concentration comes from a skewed non-clipped distribution
  (as in GCIQL, which concentrates weight on inverted transitions that never reach the ceiling).
  High fraction clipped means the useful-alpha window is narrow because many samples saturate.

  \item \textbf{Top-5\% weight mass}.
  Fraction of total weight held by the top 5\% of samples by weight.
  A direct measure of how much the policy update is dominated by a small elite set of transitions,
  regardless of whether they are clipped.

  \item \textbf{Weight entropy}.
  $-\sum_i p_i \log p_i$ where $p_i = w_i / \sum_j w_j$.
  A full-distribution summary of concentration. Higher entropy means weights are spread broadly;
  lower entropy means they are concentrated. Equivalent to $\log(\ESS)$ in the uniform-weight limit.

  \item \textbf{Alpha-sensitivity curve}.
  ESS and saturation mass recomputed at 60 log-spaced $\alpha$ values ($0.01$--$20$) for fixed advantages.
  Shows how quickly the weight distribution changes as the extraction temperature varies.
  A steep ESS collapse over a small $\alpha$ interval means the useful extraction window is narrow,
  which directly predicts a narrow trainability landscape in the $(\eta, \alpha)$ plane.

\end{itemize}

\begin{table}[htb]
\centering
\small
\renewcommand{\arraystretch}{1.1}
\begin{tabular}{lccccc}
\toprule
Algorithm & $\rho_{0.9}$ all & $\rho_{0.8}$ all & Max $\Rdist$ & $\rho_{0.9}$ last & $\rho_{0.8}$ last \\
\midrule
CRL   & $0.51\pm0.31$ & $0.63\pm0.34$ & $0.84$ & $0.69$ & $0.80$ \\
GCIQL & $0.04\pm0.02$ & $0.07\pm0.05$ & $0.71$ & $0.05$ & $0.09$ \\
GCIVL & $0.11\pm0.09$ & $0.50\pm0.20$ & $0.78$ & $0.02$ & $0.48$ \\
QRL   & $0.52\pm0.23$ & $0.62\pm0.27$ & $0.84$ & $0.70$ & $0.77$ \\
\bottomrule
\end{tabular}
\caption{AntMaze Medium landscape context. QRL and CRL have broad near-optimal regions.}
\end{table}

\section{Advantage and Concentration Results}

\begin{table}[htb]
\centering
\small
\renewcommand{\arraystretch}{1.1}
\begin{tabular}{lcccc}
\toprule
Env Family & Alg & FR-AUC & Gap mean & EI \\
\midrule
AntMaze & QRL   & $0.55$--$0.58$ & $+0.36$--$+0.68$ & $+0.19$--$+0.25$ \\
AntMaze & GCIVL & $0.56$--$0.58$ & $+0.14$ & $+0.12$--$+0.14$ \\
AntMaze & CRL   & $\approx0.53$ & $+0.17$ & $\approx+0.05$ \\
AntMaze & GCIQL & $\approx0.47$ & $-0.10$ & $-0.07$--$-0.08$ \\
\midrule
Cube & GCIVL & $0.644$ & $+0.160$ & $+0.125$ \\
Cube & CRL   & $0.562$ & $+0.875$ & $+0.164$ \\
Cube & GCIQL & $0.531$ & $-0.039$ & $-0.028$ \\
Cube & QRL   & $0.523$ & $+0.167$ & $+0.108$ \\
\bottomrule
\end{tabular}
\caption{Future-random diagnostics across environments.}
\end{table}

\begin{table}[htb]
\centering
\small
\renewcommand{\arraystretch}{1.1}
\begin{tabular}{lcccc}
\toprule
Environment & CRL & GCIVL & QRL & GCIQL \\
\midrule
AntMaze Med & $0.329$ & $0.468$ & $0.543$ & $0.140$ \\
AntMaze Lrg & $0.327$ & $0.459$ & $0.601$ & $0.144$ \\
Cube        & $0.429$ & $0.559$ & $0.470$ & $0.159$ \\
\bottomrule
\end{tabular}
\caption{Configured-alpha ESS. GCIQL is consistently the most concentrated.}
\end{table}

\begin{table}[htb]
\centering
\small
\renewcommand{\arraystretch}{1.1}
\begin{tabular}{lcccc}
\toprule
Metric (AntMaze Med) & CRL & GCIQL & GCIVL & QRL \\
\midrule
Fraction Clipped & $0.208$ & $0.089$ & $0.370$ & $0.442$ \\
Top-5\% Weight Mass & $0.236$ & $0.466$ & $0.142$ & $0.105$ \\
ESS $\leftrightarrow$ Return Corr & $-0.54^{***}$ & $+0.21^{***}$ & $+0.21^{***}$ & $+0.01$ \\
\bottomrule
\end{tabular}
\caption{Weight distribution and return correlations. Low ESS helps CRL but hurts GCIQL.}
\end{table}

\section{Interpretation and External Context}

The raw OGBench-style performance (Table \ref{tab:bench}) shows that navigation favors reachability methods (CRL/QRL), while manipulation favors value-progress (GCIVL/GCIQL).

\begin{table}[htb]
\centering
\small
\label{tab:bench}
\begin{tabular}{lcccccc}
\toprule
Dataset & GCBC & GCIVL & GCIQL & QRL & CRL & HIQL \\
\midrule
antmaze-med & $29\pm4$ & $72\pm8$ & $71\pm4$ & $88\pm3$ & $95\pm1$ & $96\pm1$ \\
antmaze-lrg & $24\pm2$ & $16\pm5$ & $34\pm4$ & $75\pm6$ & $83\pm4$ & $91\pm2$ \\
cube-single & $6\pm2$  & $53\pm4$ & $68\pm6$ & $5\pm1$  & $19\pm2$ & $15\pm3$ \\
\bottomrule
\end{tabular}
\caption{Selected raw benchmark context.}
\end{table}

On \textbf{AntMaze}, QRL/GCIVL represent \textit{diffuse extraction} of aligned progress signals. CRL performs \textit{selective extraction} where high-weight samples are reliable. GCIQL suffers from \textit{brittle concentration} on an inverted $Q-V$ tail.

On \textbf{Cube}, GCIVL is the most robust due to value-progress capturing delayed manipulation. QRL weakens as geometric alignment fails. CRL becomes \textit{scale-sensitive}: large advantages make the "safe" alpha window narrow, explaining its narrower landscape.

\section{Conclusion}

\begin{table}[t!]
\centering
\small
\begin{tabularx}{\textwidth}{lXX}
\toprule
Method & AntMaze regime & Cube regime \\
\midrule
QRL & Strong quasimetric progress; broad diffuse extraction. & Progress weakens; narrower conservative extraction. \\
GCIVL & Useful value-progress; broad/moderate diffuse extraction. & Strongest broad extraction; remains robust. \\
CRL & Selective reachability; low ESS is useful. & Sharp signal but tight alpha window; scale-sensitive. \\
GCIQL & Poorly aligned $Q-V$ tail; brittle concentration. & Inversion weakens; broadness from diluted extraction. \\
\bottomrule
\end{tabularx}
\caption{Final regime taxonomy.}
\end{table}

The paper's central claim is that trainability in offline GCRL is governed by the interaction between advantage semantics, task structure, and AWR extraction conditioning. Together, these diagnostics reveal distinct policy-extraction regimes that are invisible from final return alone.

\end{document}
